\chapter{Distributions}
\label{ch:dist}

\begin{quote}
\itshape
An object may be described by what it does to probes rather than by what value it takes at each
point. Once that shift is made, $\delta$ stops being a paradox: it is a perfectly ordinary element of
a space whose definition requires nothing more than smoothness and compact support of the probes.
\end{quote}

\section{Test functions}

\begin{definition}\label{def:D}
The space of test functions is
\begin{equation}
\D(\R)=C_{c}^{\infty}(\R)=\{\varphi:\R\to\C \;:\; \varphi\in C^{\infty},\ \supp\varphi
\text{ compact}\}.
\end{equation}
\end{definition}

The two requirements do different work. Smoothness allows derivatives of arbitrary order to be
transferred from the object to the probe, which is how differentiation will be defined. Compact
support makes boundary terms vanish in every integration by parts, and it confines the observation to
a bounded region, which is what allows $\Lonloc$ functions to act.

\begin{example}[a smooth function with compact support]\label{ex:bump}
\begin{equation}
\eta(x)=\begin{cases}\exp\!\left(-\dfrac{1}{1-x^{2}}\right),&|x|<1,\\[8pt] 0,&|x|\ge1 .\end{cases}
\end{equation}
As $x\to\pm1$ the exponent tends to $-\infty$ faster than any power, so $\eta$ and all its
derivatives tend to $0$ at the boundary, and the piecewise definition is $C^{\infty}$. The function
is positive on $(-1,1)$ and vanishes identically outside.
\end{example}

\begin{example}
$e^{-x^{2}}$ is smooth and decays rapidly, but its support is all of $\R$; it belongs to the Schwartz
space of Chapter~\ref{ch:tespered} and not to $\D$.
\end{example}

\begin{definition}[convergence in $\D$]
$\varphi_{j}\to\varphi$ in $\D$ means: there is a compact $K$ containing every $\supp\varphi_{j}$ and
$\supp\varphi$, and for every $k\ge0$,
\begin{equation}
\sup_{x\in K}\big|\varphi_{j}^{(k)}(x)-\varphi^{(k)}(x)\big|\longrightarrow0 .
\end{equation}
\end{definition}

The common compact set is part of the definition. A sequence of probes that shrinks while moving off
to infinity is not convergent in $\D$, and admitting such sequences would destroy the continuity
properties the theory needs.

\section{Distributions}

\begin{definition}\label{def:distribution}
A \emph{distribution} on $\R$ is a map $T:\D(\R)\to\C$, written $\varphi\mapsto\ip{T}{\varphi}$,
which is
\begin{enumerate}
\item \emph{linear}: $\ip{T}{a\varphi+b\psi}=a\ip{T}{\varphi}+b\ip{T}{\psi}$ for $a,b\in\C$;
\item \emph{continuous}: $\varphi_{j}\to\varphi$ in $\D$ implies $\ip{T}{\varphi_{j}}\to\ip{T}{\varphi}$.
\end{enumerate}
The space of distributions is denoted $\D'(\R)$.
\end{definition}

\begin{remark}[the pairing is bilinear, not Hermitian]
There is no complex conjugation in $\ip{T}{\varphi}$. It is linear in each slot over $\C$. In
particular $\ip{T}{\varphi}$ is not an inner product, and statements such as ``$\F$ is unitary'' have
no meaning in this framework. The absence of conjugation is what fixes the position of the factors
$2\pi$ in Chapter~\ref{ch:tespered}.
\end{remark}

\begin{proposition}[local estimate]\label{prop:localest}
A linear $T:\D\to\C$ is a distribution if and only if for every compact $K$ there exist $C_{K}<\infty$
and $m_{K}\in\mathbb{N}$ such that
\begin{equation}
|\ip{T}{\varphi}|\le C_{K}\max_{0\le k\le m_{K}}\|\varphi^{(k)}\|_{\infty}
\qquad\text{whenever }\supp\varphi\subset K .
\label{eq:localest}
\end{equation}
\end{proposition}

\begin{proof}[Sketch]
Sufficiency: \eqref{eq:localest} makes continuity immediate, since convergence in $\D$ forces every
derivative to converge uniformly on $K$. Necessity: if no such estimate existed, one could construct
$\varphi_{j}\to0$ in $\D$ with $|\ip{T}{\varphi_{j}}|\ge1$, contradicting continuity.
\end{proof}

The estimate expresses the only kind of control required: on a fixed region, the response is bounded
by finitely many derivatives of the probe. The integer $m_{K}$ is the \emph{order} of $T$ on $K$.

\section{Ordinary functions as distributions}

\begin{theorem}[embedding]\label{thm:embedding}
Let $f\in\Lonloc$. Then
\begin{equation}
\ip{T_{f}}{\varphi}:=\int_{\R}f(x)\varphi(x)\dd x
\end{equation}
defines a distribution. Moreover $T_{f}=T_{g}$ if and only if $f=g$ almost everywhere.
\end{theorem}

\begin{proof}
Since $\varphi$ is supported in a compact $K$, the integral is over $K$ and is finite by local
integrability. The estimate
\begin{equation}
|\ip{T_{f}}{\varphi}|\le\|f\|_{L^{1}(K)}\|\varphi\|_{\infty}
\end{equation}
is \eqref{eq:localest} with $m_{K}=0$. For the last statement, if $\int(f-g)\varphi=0$ for all
$\varphi\in\D$ supported in $K$, then $f-g=0$ almost everywhere on $K$; this is the standard
duality between locally integrable functions and test functions. Taking a sequence of compact sets
exhausting $\R$ gives the result.
\end{proof}

\begin{remark}[values at a point are invisible]
The theorem says that a distribution sees only the almost-everywhere class of an $\Lonloc$ function.
In particular the value $H(0)$ is irrelevant: changing it changes the function on a null set, hence
not the distribution. This is not vagueness; it is a precise statement, and it removes a question
that would otherwise be unanswerable.
\end{remark}

\section{\texorpdfstring{$\delta$}{delta} is not a function}

\begin{theorem}\label{thm:delta-not-L1}
There is no $g\in\Lonloc$ with $T_{g}=\delta$.
\end{theorem}

\begin{proof}
Suppose $T_{g}=\delta$. Let $\varphi\in C_{c}^{\infty}$ be supported in $\R\setminus\{0\}$. Then
$\ip{\delta}{\varphi}=0$, so $\int g\varphi=0$. As $\varphi$ ranges over test functions supported in
$\R\setminus\{0\}$, Theorem~\ref{thm:embedding} gives $g=0$ almost everywhere on $\R\setminus\{0\}$.
Since $\{0\}$ is a null set, $g=0$ almost everywhere on $\R$, hence $T_{g}=0\neq\delta$, a
contradiction.
\end{proof}

So $\delta$ is a genuinely new object: it is not a function with a strange value at one point, it is
not any function at all. Two remarks prevent over-reading this.

\begin{enumerate}
\item ``Not a function'' does not mean ``not defined''. $\delta$ is defined by
\eqref{eq:deltadef}, and it is also a measure: the Dirac point mass, assigning $1$ to a set
containing $0$ and $0$ otherwise. Measures embed into distributions, and $\delta$ is the simplest
member of the intersection.
\item The failure in Theorem~\ref{thm:delta-not-L1} is a failure of \emph{representation}, not of
existence. Every $\Lonloc$ function gives a distribution; $\delta$ is a distribution that no
$\Lonloc$ function gives.
\end{enumerate}

\section{Operations}

\subsection*{Differentiation}

\begin{definition}
$\ip{T'}{\varphi}:=-\ip{T}{\varphi'}$, and $\ip{T^{(n)}}{\varphi}:=(-1)^{n}\ip{T}{\varphi^{(n)}}$.
\end{definition}

This is well defined because $\varphi\in\D$ implies $\varphi'\in\D$, and because
$\varphi_{j}\to\varphi$ in $\D$ implies $\varphi_{j}'\to\varphi'$ in $\D$.

\begin{proposition}[differentiation preserves convergence]\label{prop:derivcont}
If $T_{j}\to T$ in $\D'$, then $T_{j}'\to T'$ in $\D'$.
\end{proposition}

\begin{proof}
$\ip{T_{j}'}{\varphi}=-\ip{T_{j}}{\varphi'}\to-\ip{T}{\varphi'}=\ip{T'}{\varphi}$ for every
$\varphi$.
\end{proof}

\subsection*{Multiplication by a smooth function}

For $a\in C^{\infty}(\R)$ define
\begin{equation}
\ip{aT}{\varphi}:=\ip{T}{a\varphi},
\end{equation}
which is well defined because $a\varphi\in\D$. The operation is associative and compatible with
differentiation.

\begin{proposition}\label{prop:mult-rules}
For $a\in C^{\infty}$ and any distribution $T$,
\begin{align}
a\,\delta&=a(0)\,\delta, \label{eq:adelta}\\
x\,\delta&=0, \label{eq:xdelta}\\
x\,\delta'&=-\delta, \label{eq:xdelta'}\\
x\,\delta^{(n)}&=-n\,\delta^{(n-1)}, \label{eq:xdeltaN}\\
(aT)'&=a'T+aT' . \label{eq:leibniz}
\end{align}
\end{proposition}

\begin{proof}
For \eqref{eq:adelta}: $\ip{a\delta}{\varphi}=\ip{\delta}{a\varphi}=a(0)\varphi(0)=a(0)\ip{\delta}{\varphi}$.

For \eqref{eq:xdelta'}: $\ip{x\delta'}{\varphi}=\ip{\delta'}{x\varphi}=-(x\varphi)'(0)=-\varphi(0)$.

For \eqref{eq:xdeltaN}: since $(x\varphi)^{(n)}=x\varphi^{(n)}+n\varphi^{(n-1)}$, evaluating at $0$
leaves $n\varphi^{(n-1)}(0)$, so
\begin{equation}
\ip{x\delta^{(n)}}{\varphi}=(-1)^{n}n\varphi^{(n-1)}(0)
=-n\,(-1)^{n-1}\varphi^{(n-1)}(0)=-n\ip{\delta^{(n-1)}}{\varphi}.
\end{equation}

For \eqref{eq:leibniz}: using $(a\varphi)'=a'\varphi+a\varphi'$,
\begin{equation}
\ip{(aT)'}{\varphi}=-\ip{T}{(a\varphi)'}
=-\ip{T}{a'\varphi}-\ip{T}{a\varphi'}
=\ip{a'T}{\varphi}+\ip{aT'}{\varphi}. \qedhere
\end{equation}
\end{proof}

\begin{remark}
Formula \eqref{eq:xdeltaN} explains Proposition~\ref{prop:keromega} a second time: multiplying by
$\omega$ lowers the order of $\delta^{(n)}$ by one, so among the $\delta^{(n)}$ only $\delta$ itself
is annihilated.
\end{remark}

\subsection*{Support}

Definition~\ref{def:support} applies verbatim to distributions, with $\supp\varphi\subset U$. The
support of $\delta^{(n)}$ is $\{0\}$ (Proposition~\ref{prop:suppdelta}); the support of $H$ is
$[0,\infty)$; the support of a polynomial is all of $\R$ as a function, but its \emph{spectrum} is
supported at the origin, and the two uses of the word must not be confused.

\section{Translation, scaling, composition}

\begin{proposition}\label{prop:affine}
Let $a\neq0$ and $b\in\R$. Then
\begin{equation}
\ip{\delta_{b}}{\varphi}=\varphi(b),
\qquad
\delta(ax)=\frac{1}{|a|}\,\delta(x),
\qquad
\delta(x-b)=\delta_{b}(x).
\label{eq:delta-scale}
\end{equation}
\end{proposition}

\begin{proof}
For the scaling, test against $\varphi$ and substitute $u=ax$:
\begin{equation}
\ip{\delta(a\cdot)}{\varphi}=\left\langle\delta,\frac{1}{|a|}\varphi\!\left(\frac{\cdot}{a}\right)\right\rangle
=\frac{1}{|a|}\varphi(0)=\left\langle\frac{1}{|a|}\delta,\varphi\right\rangle . \qedhere
\end{equation}
\end{proof}

The factor $|a|^{-1}$ is a Jacobian, not an algebraic reciprocal; the absolute value matters when
$a<0$.

\begin{proposition}[composition]\label{prop:composition}
Let $g\in C^{\infty}$ have only simple zeros $x_{j}$, locally finite in number. Then
\begin{equation}
\delta\big(g(x)\big)=\sum_{j}\frac{\delta(x-x_{j})}{|g'(x_{j})|} .
\label{eq:compose}
\end{equation}
\end{proposition}

\begin{proof}[Sketch]
Near a simple zero $x_{j}$, write $g(x)=(x-x_{j})h_{j}(x)$ with $h_{j}(x_{j})=g'(x_{j})\neq0$; the
change of variable $u=g(x)$ on a small interval around $x_{j}$ has Jacobian
$1/|g'(x_{j})|$, and outside the union of these intervals the test function is annihilated because
$g$ does not vanish there. Summing the contributions gives \eqref{eq:compose}.
\end{proof}

\begin{remark}
The hypothesis of simple zeros cannot be dropped. For $g(x)=x^{2}$ the derivative vanishes at the
zero, the formula divides by zero, and indeed $\delta(x^{2})$ is not of this form.
\end{remark}

\section{Operations that are not available}

\begin{enumerate}
\item \textbf{Products of distributions.} There is no general definition of $ST$ for
$S,T\in\D'$. In particular $\delta^{2}$ has no canonical meaning, and expressions such as
$\|\delta\|^{2}$ or ``the energy of $\delta$'' are meaningless.
\item \textbf{Multiplication by nonsmooth functions.} $H\delta$ has no intrinsic definition: $H$ is
not smooth, and different regularizations of the product give different answers. Choosing $H(0)$
does not resolve this.
\item \textbf{Convolutions of arbitrary distributions.} A compactly supported distribution may be
convolved with any distribution; two arbitrary distributions need not admit a convolution.
\item \textbf{Confusion with measures.} Distributions are not signed measures. $\delta$ is both, but
$\delta'$ is not a measure, and neither is $\PV\frac{1}{\omega}$ in the sense of positive or signed
measure with finite local variation (Proposition~\ref{prop:Kprime}(iv) makes this visible: the
total variation diverges).
\end{enumerate}

\section{Regularization: every distribution is a limit of smooth functions}

One structural fact remains, and it explains why the method of \emph{approximating by ordinary
functions} never loses information.

Let $\rho\in C_{c}^{\infty}$ with $\rho\ge0$, $\int\rho=1$, and set
$\rho_{\ep}(x)=\ep^{-1}\rho(x/\ep)$.

\begin{theorem}\label{thm:regularization}
Let $T$ be a distribution with compact support. For $\ep>0$ define
\begin{equation}
T_{\ep}(x):=\big\langle T,\;\rho_{\ep}(x-\cdot)\big\rangle .
\label{eq:Teps}
\end{equation}
Then $T_{\ep}\in C^{\infty}(\R)$, and $T_{\ep}\to T$ in $\D'(\R)$ as $\ep\downarrow0$.
\end{theorem}

\begin{proof}
\emph{Smoothness.} Let $K$ be a compact set containing $\supp T$ in its interior. For $x$ in a
bounded set, the map $y\mapsto\rho_{\ep}(x-y)$ is supported in a fixed compact neighbourhood of $K$, so
the pairing is legitimate. Differentiation under the pairing is allowed because
$\partial_{x}^{\alpha}\rho_{\ep}(x-y)=(-\partial_{y})^{\alpha}\rho_{\ep}(x-y)$ is again a test
function, and $T$ is continuous; hence $T_{\ep}\in C^{\infty}$.

\emph{Convergence.} Let $\varphi\in C_{c}^{\infty}$. Interchanging the pairing with the integration in
$x$ (a finite sum after expanding on a compact set),
\begin{equation}
\ip{T_{\ep}}{\varphi}
=\int_{\R}\big\langle T,\rho_{\ep}(x-\cdot)\big\rangle\varphi(x)\dd x
=\Big\langle T,\;\int_{\R}\rho_{\ep}(x-\cdot)\varphi(x)\dd x\Big\rangle
=\big\langle T,\;\rho_{\ep}*\varphi\big\rangle ,
\end{equation}
where the last step uses the symmetry $\rho_{\ep}(-y)=\rho_{\ep}(y)$. By
Theorem~\ref{thm:conv}, $\rho_{\ep}*\varphi\to\varphi$ in $C^{\infty}$ uniformly on compact sets, hence
in $\D$, and continuity of $T$ gives $\ip{T_{\ep}}{\varphi}\to\ip{T}{\varphi}$.
\end{proof}

\begin{remark}
The hypothesis of compact support is what allows the mollification to be tested against $T$ without
worrying about behaviour at infinity. For general $T$ a localization argument---multiply by a smooth
cutoff equal to one on a large set---reduces to this case, at the cost of a statement that holds on
compact subsets only, which is exactly what convergence in $\D'$ means.
\end{remark}

\begin{corollary}\label{cor:approx}
Every result in this book that was obtained by regularizing with a specific family
($e^{-\ep|t|}$, $H e^{-\ep t}$, $g_{\sigma}$, $w_{T}$, $K_{\ep}$) is an instance of a single
mechanism: convolving, or pairing, against a mollifier.
\end{corollary}

\section{Summary of the chapter}

\begin{center}
\small
\begin{tabular}{@{}p{3.4cm}p{9.4cm}@{}}
\toprule
Notion & Content \\
\midrule
Test functions & $C_{c}^{\infty}$: smooth (to move derivatives) and compactly supported (to kill
boundary terms) \\
Distribution & continuous linear functional on $\D$; continuity estimated by
\eqref{eq:localest} \\
Embedding & $\Lonloc\hookrightarrow\D'$, injective modulo null sets
(Theorem~\ref{thm:embedding}) \\
$\delta$ & a distribution, not an $\Lonloc$ function
(Theorem~\ref{thm:delta-not-L1}); also a measure \\
Derivative & transfer to the probe with $(-1)^{n}$; continuous for distributional convergence \\
Multiplication & only by smooth functions; \eqref{eq:adelta}--\eqref{eq:leibniz} \\
Forbidden & products of distributions, $\delta^{2}$, $H\delta$, unrestricted convolution \\
\bottomrule
\end{tabular}
\end{center}
